% Abhakseite „Das kann ich“ – Zone nur im Gesamt (\mitzone)
%% TODO Ich-kann-Sätze: Platzhalter wie in den Aufgabendateien
\begin{abhakseite}
\ifdefined\mitzone
\abhakgruppe{Kennst du schon}
\abhak{1}{Punkte im räumlichen Koordinatensystem lesen und eintragen, Koordinatenebenen und Achsen benennen, ein Schrägbild lesen (welche Kante liegt in welcher Koordinatenebene, welche Koordinate ist auf einer Seitenfläche konstant)}
\abhak{2}{Verbindungsvektor zweier Punkte bilden (Spitze minus Fuß), Vektoren addieren und mit einer Zahl vervielfachen, Ortsvektor und Punkt unterscheiden}
\abhak{3}{Kollinearität zweier Vektoren prüfen (ist der eine ein Vielfaches des anderen?)}
\abhak{4}{Skalarprodukt zweier Vektoren in Koordinaten berechnen und Orthogonalität als „Skalarprodukt null“ erkennen}
\abhak{5}{Geradengleichung in Parameterform lesen (Stütz- und Richtungsvektor) und die Punktprobe an einer Geraden}
\abhak{6}{Lineare Gleichungssysteme mit zwei und drei Unbekannten lösen, auch unterbestimmte (eine Unbekannte frei wählen, Vielfache als gleichwertige Lösungen erkennen)}
\abhak{7}{Eine lineare Gleichung als Punktmenge lesen}
\abhak{8}{Verbindungsvektor zweier Punkte bilden (Spitze minus Fuß), Vektoren addieren und mit einer Zahl vervielfachen, Ortsvektor und Punkt unterscheiden – Fehler finden}
\abhak{9}{Verbindungsvektor zweier Punkte bilden (Spitze minus Fuß), Vektoren addieren und mit einer Zahl vervielfachen, Ortsvektor und Punkt unterscheiden – selbst rechnen}
\fi
\abhakgruppe{Einheit 1 · Parameterform einer Ebene}
\abhak{10}{Parameterform}
\abhak{11}{Parameterform – Fehler finden, Begründen, Anwendung, Darstellungswechsel}
\abhakgruppe{Einheit 2 · Normalenvektor und Koordinatengleichung}
\abhak{12}{Koordinatengleichung bestimmen}
\abhak{13}{Normalenvektor und Nachweise}
\abhak{14}{Koordinatengleichung bestimmen – Fehler finden, Begründen, Anwendung, Darstellungswechsel}
\abhakgruppe{Einheit 3 · Ebenen im Koordinatensystem}
\abhak{15}{Besondere Lagen}
\abhak{16}{Besondere Lagen – Fehler finden, Begründen, Anwendung, Darstellungswechsel}
\abhakgruppe{Einheit 4 · Parallele Ebenen}
\abhak{17}{Parallele Ebenen}
\abhak{18}{Parallele Ebenen über die Normalenvektoren auswählen und über die Konstante der Koordinatengleichung der Höhe nach zuordnen}
\abhak{19}{Parallele Ebenen – Fehler finden, Begründen, Anwendung}
\end{abhakseite}
